\chapter{What gets installed where}
\label{ch:i-overview}

An X30 patrol system is four computers and a handheld. The robot carries three: the
\textbf{motion host} walks the robot and is never changed; the \textbf{navigation host} on the
robot's back runs the stack that drives it; the \textbf{perception host} inside the body runs
the inspection cameras and records them. The operator's \textbf{Mac or laptop} runs NETRA, the
command center, and the \textbf{Skydroid H12} runs the X30Controller app. Everything meets on the
robot's own network: its Wi-Fi, or a cable to the Debug Ethernet port at the tail.

\begin{figure}[h]
  \centering
  \includegraphics[width=\linewidth]{network}
  \caption{The computers of an X30 deployment and what talks to what.}
  \label{fig:i-network}
\end{figure}

\begin{center}
\begin{tabularx}{\linewidth}{@{}>{\bfseries}p{0.24\linewidth}p{0.22\linewidth}X@{}}
  \toprule
  Computer & Address & What you install \\
  \midrule
  Navigation host & \code{192.168.1.106} & x30\_ros (the whole workspace), the service
    \code{x30-navigation}; Deep Robotics' own autonomy turned off \\
  Perception host & \code{192.168.1.105} & x30\_ros (the camera drivers and the recorder only),
    the service \code{x30-perception} \\
  Motion host & \code{192.168.1.103} & nothing: one line in its \file{network.toml} \\
  Mac or laptop & the robot's network & NETRA (Docker), the tools to pack the sources and to
    install the app \\
  Skydroid H12 & the robot's Wi-Fi & the X30Controller app (APK) \\
  \bottomrule
\end{tabularx}
\end{center}

\section{The versions this guide installs}

\begin{center}
\begin{tabularx}{\linewidth}{@{}>{\bfseries}p{0.24\linewidth}p{0.14\linewidth}X@{}}
  \toprule
  Part & Version & Where it comes from \\
  \midrule
  x30\_ros & 0.19 & \url{https://github.com/DarmawanRobotics/x30_ros} \\
  X30Controller & 1.0 & \url{https://github.com/DarmawanRobotics/X30Controller/releases} \\
  NETRA & with the substation module & \url{https://github.com/DarmawanWeb/Netra} (private) \\
  Robot & Jueying X30 Pro & Deep Robotics; Ubuntu 20.04 and ROS Noetic on both hosts \\
  \bottomrule
\end{tabularx}
\end{center}

\section{The way through this guide}

The robot has no internet. Everything that downloads is done first, on the Mac, while it is
online; then the Mac joins the robot's Wi-Fi and carries it over.

\begin{enumerate}
  \item \textbf{Before you start} (Chapter~\ref{ch:i-before}): the parts, the addresses and the
    passwords.
  \item \textbf{On the Mac, online} (Chapter~\ref{ch:i-mac}): pack the robot's sources, download
    the app, build NETRA and register the robot.
  \item \textbf{On the robot} (Chapter~\ref{ch:i-robot}): copy, build, turn Deep Robotics'
    autonomy off, introduce the hosts to each other.
  \item \textbf{The services} (Chapter~\ref{ch:i-services}): the stack starts at boot under a
    watchdog; its settings.
  \item \textbf{The handheld} (Chapter~\ref{ch:i-handheld}) and \textbf{NETRA}
    (Chapter~\ref{ch:i-netra}).
  \item \textbf{A new robot} (Chapter~\ref{ch:i-bringup}): the checks before it walks on its own,
    the first site and route.
\end{enumerate}

\begin{dangerbox}[Safety]
The robot walks among energized equipment. Do the first bring-up in an open area, away from
live parts, with one person at the handheld whose only job is to stop the robot. Keep the
robot's hard emergency stop (at the tail) and the handheld's STOP within reach at every step.
\end{dangerbox}
